%% ch5_IV  第四章（续）：紧化上的截断与水平结构，原书 pp. 27--33
% ============================================================
%  ch5_IV_b  第 IV 章续（截断、水平结构与稳定化），及第 V 章开头
%  原书 pp. 28--36，即 PDF 第 29--37 页
%  底稿：transcript/partD.md
% ============================================================

\origpage{28}\begin{lemma*}
(i) 下列两个条件\footnote{译注：第二个条件里的分母原文印作 $s$，而首条公式里的分母为 $r$；此处照录原文。[来源：转录稿 p.~28 的「不确定清单」]}：
\[
\left\{
\begin{array}{l}
\bullet\ \text{对每个 } s\in\underline{r}^{-}\cup\underline{r}^{+},\ \text{有}\\[6pt]
\qquad p(s)-1<\deg\mathcal{E}_{s}-\dfrac{s}{r}\deg\mathcal{E}\le p(s).\\[10pt]
\bullet\ \text{对 }\tilde{\mathcal{E}}\text{ 的每个「好」子对象 }(\mathcal{F},\mathcal{F}')\text{，有}\\[6pt]
\qquad \deg\mathcal{F}-\dfrac{\mathrm{rg}\,\mathcal{F}}{s}\deg\mathcal{E}\le p(\mathrm{rg}\,\mathcal{F})\\[6pt]
\text{若 }(\mathcal{F},\mathcal{F}')\text{ 为类型 I，且}\\[10pt]
\qquad \deg\mathcal{F}-\dfrac{\mathrm{rg}\,\mathcal{F}}{s}\deg\mathcal{E}\le p(\mathrm{rg}\,\mathcal{F})-1\\[6pt]
\text{若 }(\mathcal{F},\mathcal{F}')\text{ 为类型 II。}
\end{array}
\right.
\]
在 $\mathrm{Cht}^{r}_{\underline{r}}$ 中定义出一个开集 $\mathrm{Cht}^{r,\bar{p}\le p}_{\underline{r}}$。

(ii) 存在 $\overline{\mathrm{Cht}^{r}}$ 的唯一开集 $\overline{\mathrm{Cht}^{r,\bar{p}\le p}}$，使它在每个 $\mathrm{Cht}^{r}_{\underline{r}}$ 中的迹恰为 $\mathrm{Cht}^{r,\bar{p}\le p}_{\underline{r}}$。
\end{lemma*}

\begin{remark*}
特别地，$\overline{\mathrm{Cht}^{r,\bar{p}\le p}}\cap\mathrm{Cht}^{r}=\mathrm{Cht}^{r,\bar{p}\le p}$。
\end{remark*}

我们给出各分层上截断的另一种描述。

对每个整数 $r'$，都有可表、有限、满且根式的态射\footnote{译注：原文把两个方块图画在同一张图里：左图为子丛 $\mathcal{E}'\hookrightarrow\mathcal{E}$ 连同指向 $\tau\mathcal{E}$ 的对角箭头，右图为 $\mathcal{E}'\hookrightarrow\tau\mathcal{E}$ 连同对角箭头 $\tau\mathcal{E}'\to\tau\mathcal{E}$。转录稿以 array 环境复原。[来源：转录稿图注 p.~28]}
\[
{}^{r'}\mathrm{Cht}\to\mathrm{Cht}^{r'},\qquad
\left(\begin{array}{ccc}
\mathcal{E}' & \hookrightarrow & \mathcal{E}\\
 & \searrow & \\
 & & \tau\mathcal{E}
\end{array}\right)
\longmapsto
\left(\begin{array}{ccc}
\mathcal{E}' & \hookrightarrow & \tau\mathcal{E}\\
 & & \uparrow\\
 & \tau\mathcal{E}' &
\end{array}\right)
\]
因此有有限、满且根式的态射
\[
\begin{aligned}
\mathrm{Cht}^{r}_{\underline{r}} &\to \mathrm{Cht}^{r_{1}}\times_{X}{}^{r_{2}}\mathrm{Cht}\times_{X}\cdots\times_{X}{}^{r_{k}}\mathrm{Cht}\\
&\to \mathrm{Cht}^{r_{1}}\times_{X}\mathrm{Cht}^{r_{2}}\times_{X,\mathrm{Frob}}\cdots\times_{X,\mathrm{Frob}}\mathrm{Cht}^{r_{k}}
\end{aligned}
\]
它位于 $\mathrm{Id}_{X}\times\mathrm{Frob}_{X}$ 之上。

回到截断多边形 $p:[0,r]\to\mathbb{R}_{+}$，并固定一个次数 $d\in\mathbb{Z}$。

\origpage{29}对 $0\le s\le r$，设 $d(s)$ 为满足
\[
d(s)-\frac{s}{r}d\in\,]p(s)-1,p(s)]
\]
的唯一整数。

令\footnote{译注：原文第一行印作 $d_{1}=d(r_{1})=d(r_{1})-d(0)$。因 $d(0)=0$，此式只是恒等式；疑原文有笔误，此处照录。[来源：转录稿 p.~29]}
\[
\begin{aligned}
d_{1} &= d(r_{1}) = d(r_{1})-d(0)\\
p_{1}(0) &= p_{1}(r_{1}) = 0\\
p_{1}(s) &= d(s)-\frac{s}{r_{1}}d(r_{1}),\quad 1\le s\le r_{1},
\end{aligned}
\]
并对 $2\le i\le k$ 令
\[
\begin{aligned}
d_{i} &= d(r_{1}+\cdots+r_{i})-d(r_{1}+\cdots+r_{i-1})-1\\
p_{i}(0) &= p_{i}(r_{i}) = 0\\
p_{i}(s) &= d(r_{1}+\cdots+r_{i-1}+s)-d(r_{1}+\cdots+r_{i-1})-1\\
&\quad -\frac{s}{r_{i}}\big[d(r_{1}+\cdots+r_{i})-d(r_{1}+\cdots+r_{i-1})-1\big],\quad 1\le s\le r_{i}.
\end{aligned}
\]

\begin{lemma*}
开集 $\mathrm{Cht}^{r,d,\bar{p}\le p}_{\underline{r}}\subset\mathrm{Cht}^{r,d}_{\underline{r}}$ 是开集
\[
\mathrm{Cht}^{r_{1},d_{1},\bar{p}\le p_{1}}\times\cdots\times\mathrm{Cht}^{r_{k},d_{k},\bar{p}\le p_{k}}
\]
沿有限、满且根式的态射
\[
\mathrm{Cht}^{r}_{\underline{r}}\to\mathrm{Cht}^{r_{1}}\times_{X}\mathrm{Cht}^{r_{2}}\times_{X,\mathrm{Frob}}\cdots\times_{X,\mathrm{Frob}}\mathrm{Cht}^{r_{k}}
\]
的原像。
\end{lemma*}

\begin{remark*}
这意味着 $\mathrm{Cht}^{r,d,\bar{p}\le p}_{\underline{r}}$ 的上同调是「$r$-可忽略」的。
\end{remark*}

若斜率跳跃
\[
[p(s)-p(s-1)]-[p(s)-p(s+1)]=\partial p(s)
\]
全都足够大，则称多边形 $p:[0,r]\to\mathbb{R}_{+}$ 为「足够凸」。

\begin{proposition*}
若 $p$ 足够凸（视 $X$ 的亏格而定），则态射
\origpage{30}
\[
\overline{\mathrm{Cht}^{r,\bar{p}\le p}}\to X\times X\times(\mathbb{A}^{1}/\mathbb{G}_{m})^{r-1}
\]
光滑。
\end{proposition*}

\begin{theorem*}
若 $p$ 足够凸（$\partial p(s)\ge 2$，$\forall s$），则态射
\[
\overline{\mathrm{Cht}^{r,d,\bar{p}\le p}}\to X\times X
\]
有限型且紧合（特别地，分离）。
\end{theorem*}

\begin{theorem*}
在
\[
\overline{\mathrm{Cht}^{r,\bar{p}\le p}}/a^{\mathbb{Z}}\cong\coprod_{1\le d\le r|\deg(a)|}\overline{\mathrm{Cht}^{r,d,\bar{p}\le p}},
\]
中，Hecke 对应 $f\in\mathcal{H}^{r}_{\emptyset}$\footnote{译注：原文下标印作 $\emptyset$（斜杠圆圈），也可能本意为 $0$；此处照录。[来源：转录稿 p.~30 的「不确定清单」]} 在「其不动点附近」保持 chtouca 开集
\[
\mathrm{Cht}^{r,\bar{p}\le p}/a^{\mathbb{Z}}
\]
不变。
\end{theorem*}

现在转向
\par\medskip\noindent\textbf{水平结构}\ 
设 $N\hookrightarrow X$ 为一个水平，即一个有限闭子概形。

\par\medskip\noindent\textbf{从 $X$ 到 $N$ 的限制态射}\ 
函子
\[
\mathcal{E}\,(\text{即 }X\times S\text{ 上的丛})\mapsto\mathcal{E}\otimes_{\mathcal{O}_{X\times S}}\mathcal{O}_{N\times S}=\mathcal{E}\otimes_{\mathcal{O}_{X}}\mathcal{O}_{N}=\mathcal{E}_{N}
\]
定义出态射
\[
\mathrm{Res}:\overline{\mathrm{Cht}^{r}}\times_{X\times X}(X-N)\times(X-N)\xrightarrow{\ \cdot\otimes_{\mathcal{O}_{X}}\mathcal{O}_{N}\ }\mathcal{C}^{r,N}
\]
其中 $\mathcal{C}^{r,N}$ 是分类下列各项的叠：

\begin{itemize}
\item $\mathcal{E}$：$N\times S$ 上秩 $r$ 的丛，
\item \origpage{31}$\mathcal{L}_{1},\ldots,\mathcal{L}_{r-1}$：$S$ 上的可逆层，带截面 $\ell_{1},\ldots,\ell_{r-1}$，
\item 完整同态
\[
{}^{\tau}\mathcal{E}\Rightarrow\mathcal{E}
\]
它在 $(\mathbb{A}^{1}/\mathbb{G}_{m})^{r-1}(N\times S)$ 中的像为
\[
(\mathcal{L}_{1},\ell_{1})^{\otimes(q-1)},\ldots,(\mathcal{L}_{r-1},\ell_{r-1})^{\otimes(q-1)}.
\]
\end{itemize}

具体言之，它由处处非零的同态
\[
{}^{\tau}\Big(\Lambda^{s}\mathcal{E}\otimes\bigotimes_{t<s}\mathcal{L}_{t}^{\otimes(s-t)}\Big)\xrightarrow{u_{s}}\Lambda^{s}\mathcal{E}\otimes\bigotimes_{t<s}\mathcal{L}_{t}^{\otimes(s-t)},\quad 1\le s\le r,
\]
组成，这些同态特别满足
\[
\Lambda^{s}u_{1}=\Big(\prod_{t<s}\ell_{t}^{(q-1)(s-t)}\Big)\cdot u_{s}.
\]

条件
\[
{}^{\tau}\mathcal{E}\Rightarrow\mathcal{E}\ \text{是同构 }(\Leftrightarrow\ell_{1},\ldots,\ell_{r-1}\ \text{可逆})
\]
在 $\mathcal{C}^{r,N}$ 中定义出一个稠密开集
\[
\mathcal{C}^{r,N}_{\emptyset}=\bullet/\mathrm{GL}_{r}(\mathcal{O}_{N})=\text{$\mathrm{GL}_{r}(\mathcal{O}_{N})$ 的分类叠}
\]
它以
\[
\mathcal{C}^{r}_{N,\emptyset}=\bullet=\mathrm{Spec}\,\mathbb{F}_{q}
\]
为有限 étale Galois 覆盖。

$\mathrm{Cht}^{r}$ 是 $\mathcal{C}^{r,N}_{\emptyset}$ 沿 $\mathrm{Res}$ 的原像，并且
\[
\mathrm{Cht}^{r}_{N}=[\mathrm{Cht}^{r}\times_{X\times X}(X-N)\times(X-N)]\times_{\mathcal{C}^{r,N}_{\emptyset}}\mathcal{C}^{r}_{N,\emptyset}.
\]

\begin{theorem*}
(i) 若 $p$ 足够凸（视 $X$ 与 $N$ 的亏格而定），则
\[
(q,\mathrm{Res}):(X-N)\times(X-N)\times_{X\times X}\overline{\mathrm{Cht}^{r,\bar{p}\le p}}/a^{\mathbb{Z}}\to(X-N)\times(X-N)\times\mathcal{C}^{r,N}
\]
光滑。

\origpage{32}(ii) 若 $\mathcal{L}$ 是 $\mathcal{C}^{r,N}$ 上的 $\ell$-进层，且支撑于边界 $\mathcal{C}^{r,N}-\mathcal{C}^{r,N}_{\emptyset}$，则上同调层
\[
R^{i}q_{*}(\mathrm{Res})^{*}\mathcal{L}
\]
在 $(X-N)\times(X-N)$ 上光滑，并且 $r$-可忽略。
\end{theorem*}

\begin{corollary*}
若 $\mathcal{C}^{r}_{N}$ 是 $\mathcal{C}^{r,N}$ 上的可表叠，延拓了 $\mathcal{C}^{r,N}_{\emptyset}$ 的 étale 覆盖 $\mathcal{C}^{r}_{N,\emptyset}$，则
\[
\overline{\mathrm{Cht}^{r,\bar{p}\le p}_{N}}=[(X-N)\times(X-N)\times_{X\times X}\overline{\mathrm{Cht}^{r,\bar{p}\le p}}]\times_{\mathcal{C}^{r,N}}\mathcal{C}^{r}_{N}
\]
在 $(X-N)\times(X-N)\times\mathcal{C}^{r}_{N}$ 上光滑；它以 $\mathrm{Cht}^{r,\bar{p}\le p}_{N}$ 为开集；$\mathrm{Cht}^{r,\bar{p}\le p}_{N}-\overline{\mathrm{Cht}^{r,\bar{p}\le p}_{N}}$ 的上同调 $r$-可忽略。
\end{corollary*}

若存在 $\mathcal{C}^{r}_{N}$ 使得

\begin{itemize}
\item $\mathcal{C}^{r}_{N}$ 在 $\mathcal{C}^{r,N}$ 上射影，
\item $\mathcal{C}^{r}_{N}$ 光滑，
\item $\mathcal{C}^{r}_{N}-\mathcal{C}^{r}_{N,\emptyset}=\mathrm{DCN}$\footnote{译注：$\mathrm{DCN}$ 即 diviseur à croisements normaux，正规交叉除子。原文用一个左大括号把上述三条条件括在一起，右侧分两行标注「若 $N$ 既约，则为是」与「若 $N$ 有重数，则为？」。[来源：转录稿图注 p.~32]}，
\end{itemize}

则 $\overline{\mathrm{Cht}^{r,\bar{p}\le p}_{N}}$ 满足：

\begin{itemize}
\item $\overline{\mathrm{Cht}^{r,\bar{p}\le p}_{N}}/a^{\mathbb{Z}}$ 在 $(X-N)\times(X-N)$ 上紧合，
\item 在 $(X-N)\times(X-N)$ 上光滑，
\item 边界 $=$ 相对 $\mathrm{DCN}$，各分层皆 $r$-可忽略。
\end{itemize}

设 $\mathcal{C}'{}^{r}_{N}\to\mathcal{C}^{r,N}$ 是分类下列族
\[
(\mathcal{E};\mathcal{L}_{1},\ldots,\mathcal{L}_{r-1};\ell_{1},\ldots,\ell_{r-1};{}^{\tau}\mathcal{E}\Rightarrow\mathcal{E})
\]
的叠：与上面相同，再加一个完整同态
\[
\mathcal{E}\Rightarrow\mathcal{O}^{r}_{N\times S}
\]
\origpage{33}其像在 $(\mathbb{A}^{1}/\mathbb{G}_{m})^{k-1}(N\times S)$\footnote{译注：原文此处印作 $k-1$，而上文同类公式中为 $r-1$；此处照录。[来源：转录稿 p.~33]} 中为 $(\mathcal{L}_{1},\ell_{1}),\ldots,(\mathcal{L}_{r-1},\ell_{r-1})$，并由处处非零的同态
\[
\Lambda^{s}\mathcal{E}\otimes\bigotimes_{t<s}\mathcal{L}_{t}^{\otimes(s-t)}\xrightarrow{v_{s}}\Lambda^{s}\mathcal{O}_{N\times S},\quad 1\le s\le r,
\]
表示，这些同态特别满足
\[
\Lambda^{s}v_{1}=\Big(\prod_{t<s}\ell_{t}^{(s-t)}\Big)v_{s}
\]
以及附加条件
\[
{}^{\tau}v_{s}=v_{s}\circ u_{s}.
\]

$\mathcal{C}'{}^{r}_{N}$ 映入 $\mathcal{C}^{r,N}$ 的开集 $\mathcal{C}'{}^{r,N}$，在其中各 $u_{s}$ 处处非 $\tau$-幂零。

\begin{lemma*}
$\mathcal{C}'{}^{r}_{N}$ 是 $\mathcal{C}'{}^{r,N}$ 的有限、平坦、分歧覆盖，它延拓了 $\mathcal{C}^{r,N}_{\emptyset}$ 的 étale 覆盖 $\mathcal{C}^{r}_{N,\emptyset}$。它在 $(\mathbb{A}^{1}/\mathbb{G}_{m})^{r-1}$ 上光滑。
\end{lemma*}

\begin{lemma*}
$\overline{\mathrm{Cht}^{r}}$ 中作为 $\mathcal{C}'{}^{r,N}$ 原像的开集 $\overline{\mathrm{Cht}^{r}}{}'$，也由「极点、零点与退化子都避开 $N$」这一条件定义。
\end{lemma*}

\begin{theorem*}
（稳定性）设 $\tilde{\mathcal{E}},\tilde{\mathcal{F}}$ 是 $\overline{\mathrm{Cht}^{r,\bar{p}\le p}}$ 的两个点，取值于同一个 trait\footnote{译注：trait 指离散赋值环的谱，即由泛点与闭点构成的局部概形；此处保留法文原词，待术语表统一。[来源：术语表未收]}。设它们的一般纤维 $\tilde{\mathcal{E}}_{\eta}$ 与 $\tilde{\mathcal{F}}_{\eta}$ 是两个 chtouca，互为某个 Hecke 对应的像（该对应可再复合以 $\mathrm{Frob}_{\infty}$、$\mathrm{Frob}_{0}$ 的幂）。于是，若 $x$ 是特殊化 $\tilde{\mathcal{E}}_{s}$ 的极点、零点或退化子，则存在特殊化 $\tilde{\mathcal{F}}_{s}$ 的极点、零点或退化子 $x'$，以及 $n,n'\in\mathbb{N}$，使得
\[
\mathrm{Frob}^{n}(x)=\mathrm{Frob}^{n'}(x').
\]
\end{theorem*}

\begin{corollary*}
叠
\[
\begin{aligned}
\overline{\mathrm{Cht}^{r,\bar{p}\le p}_{N}}{}'/a^{\mathbb{Z}}&=[(X-N)\times(X-N)\times_{X\times X}\overline{\mathrm{Cht}^{r,\bar{p}\le p}}/a^{\mathbb{Z}}]\times_{\mathcal{C}^{r,N}}\mathcal{C}'{}^{r}_{N}\\
&=\overline{\mathrm{Cht}^{r,\bar{p}\le p}}{}'\times_{\mathcal{C}^{r,N}}\mathcal{C}'{}^{r}_{N}
\end{aligned}
\]
满足：

\begin{itemize}
\item \origpage{34}$\mathrm{Cht}^{r,\bar{p}\le p}_{N}/a^{\mathbb{Z}}=$ 稠密开集，
\item 在 $(X-N)\times(X-N)$ 上光滑，
\item 边界 $=$ 相对 $\mathrm{DCN}$，
\item 边界（及其各分层）的上同调 $r$-可忽略，
\item Hecke 对应保持它不变，
\item 开集 $\mathrm{Cht}^{r,\bar{p}\le p}_{N}/a^{\mathbb{Z}}$ 在「不动点附近的局部」意义下保持不变。
\end{itemize}
\end{corollary*}
